\section{Related Work}
\label{sec:related_work}

\subsection{Ultrasound lesion classification}
Ultrasound lesion classification has been studied most extensively in breast and thyroid imaging, where public data availability has enabled repeated benchmarking \cite{Al_Dhabyani_2020,Zhang_2025}. Earlier breast ultrasound work often relied on handcrafted morphological and texture descriptors combined with conventional classifiers \cite{G_mez_Flores_2015}. More recent deep models have shifted the focus toward representation learning directly from lesion-centered images or regions of interest. For example, Luo \textit{et al.} proposed a multi-view learning strategy for breast ultrasound tumor classification in \textit{Pattern Recognition}, showing that task-specific view construction can improve classification accuracy \cite{Luo_2023}. These studies confirm that ultrasound lesion classification benefits from architectures that can retain subtle boundary, texture, and contextual cues. However, the majority of published systems remain strongly tied to a single organ or a single benchmark.

Compared with breast ultrasound, public thyroid and liver ultrasound classification benchmarks are less mature but increasingly important. TN5000 has substantially improved the benchmarking environment for thyroid nodule analysis by providing a large-scale public dataset and standardized evaluation protocol \cite{Zhang_2025}. The AUL dataset broadens public benchmarking to liver ultrasound by offering annotated benign, malignant, and normal cases \cite{aul_dataset}. Taken together, these datasets suggest that a useful ultrasound classification architecture should not only perform well on one organ, but also remain competitive when transferred to datasets with different anatomy, lesion appearance, and acquisition characteristics.

\subsection{Efficient visual architectures}
The design of efficient visual backbones has rapidly evolved from classical convolutional networks to compact CNN--Transformer hybrids. ResNet, DenseNet, EfficientNet, and MobileNetV3 established widely used convolutional design principles such as residual learning, dense feature reuse, compound scaling, and hardware-aware mobile optimization \cite{he2016resnet,huang2017densenet,tan2019efficientnet,howard2019mobilenetv3}. Vision Transformer and DeiT then demonstrated the strength of attention-based image recognition, while Swin Transformer and ConvNeXt clarified the practical value of hierarchical representations and convolutional inductive biases \cite{dosovitskiy2021vit,touvron2021deit,liu2021swin,liu2022convnext}. These models form the main reference space for modern visual backbones.

Recent lightweight designs focus less on raw parameter reduction and more on where computation should be allocated. MobileViT, EdgeNeXt, EfficientFormer, RepViT, and TransXNet are representative examples of this direction, combining local detail modeling, global context aggregation, and deployment-aware efficiency in different ways \cite{mehta2022mobilevit,edgenext,li2022efficientformer,repvit,transxnet}. This line of work is particularly relevant to ultrasound analysis, where lesions often occupy a small portion of the image and subtle structural changes can be diagnostically important. Our architecture is built in this efficient-backbone context rather than as a detached medical-only network.

Attention refinements are also relevant to this design space. Differential Transformer introduced differential attention as the difference between two softmax attention maps to reduce attention noise and emphasize relevant context~\cite{ye2024differential}. We do not treat this mechanism as a new primitive. Instead, we adapt it inside a compact ultrasound ROI classifier, where suppressing common speckle and background responses is useful for lesion-focused discrimination.

Medical image analysis has also begun to adopt lightweight hybrid designs. LightCM-PNet emphasizes real-time ultrasound segmentation under deployment constraints \cite{Wang_2024}, while Lite-MixedNet explores efficient CNN--Transformer fusion for medical image segmentation \cite{Ren_2025}. These works show that lightweight design is becoming a first-class consideration rather than a secondary implementation detail. Our work follows this trend, but targets ultrasound lesion classification rather than segmentation. The goal is not to propose an entirely separate paradigm, but to derive a coherent and task-motivated architecture that preserves efficiency while improving fine-grained lesion discrimination.

\subsection{Multi-dataset validation in medical image analysis}
Reporting results on a single benchmark is often insufficient in medical image analysis because performance may depend strongly on dataset-specific factors such as annotation protocol, class prevalence, acquisition device, and institution-specific image characteristics \cite{litjens2017survey,tajbakhsh2020embracing}. This issue is especially pronounced for ultrasound, where both image quality and lesion appearance can vary substantially across clinical settings. Recent medical-image-analysis studies have therefore begun to supplement their main benchmarks with robustness analysis, external validation, or deployment-oriented testing \cite{Chen_2023,Wang_2024}. Such evidence does not automatically prove universal generalization, but it gives a more credible picture of where a method is strong and where it remains fragile.

Our paper adopts this perspective. The primary claim is strong dataset-specific in-domain performance across several public ultrasound datasets, with TN5000 as the main benchmark and BUSI/AUL as additional public benchmarks on breast and liver ultrasound. BUSI and AUL help characterize performance under different anatomy, lesion appearance, and acquisition characteristics, but they are not treated as thyroid-to-breast or thyroid-to-liver external generalization tests. Temperature scaling is used only as validation-derived post-hoc calibration, following the established calibration literature \cite{guo2017calibration}. This positioning keeps the evidence chain aligned with what the experiments can defensibly support.
